\chapter{指示関数}
    \section{諸定理}
        \subsection{\texorpdfstring{$\Omega$}{Ω}を定義域とする写像\texorpdfstring{$f,g,T$}{f,g,T}に対して\texorpdfstring{$\forall x\in\Omega,\forall t\in T(\Omega),\; \indII{T(x)=t}f(x) = \indII{T(x)=t}g(x) \iff ^\forall x\in\Omega,\;f(x)=g(x)$}{∀x∈Ω, ∀t∈T(Ω), 1(T(x)=t)f(x) = 1(T(x)=t)g(x) ⇔ ∀x∈Ω, f(x)=g(x)}}
            \begin{proof}
                \quad\\
                ($\Rightarrow$)
                \par
                $x\in\Omega$を任意にとる。$t=T(x)$を満たす$t\in T(\Omega)$をとって主張の左側の式に代入すると右側の式を得る。\\
                ($\Leftarrow$)
                \par
                $x\in\Omega,t\in T(\Omega)$を任意にとる。$t=T(x)$の場合は主張の右側の式より左側の式が成り立つ。
                $t\neq T(x)$の場合は主張の左側の式は$0=0$となる。
            \end{proof}